\section*{Capítulo \ref{ch:g12:climate-change} --- El cambio climático: observaciones, causas y proyecciones}

\begin{solution}{exo:g12:climate-change:1}
El cambio que provoca en el flujo neto de energía que entra en el sistema climático
en lo alto de la atmósfera, antes de que responda la temperatura de la superficie;
se mide en \unit{W/m^2}. Positivo: el planeta gana energía y se calienta;
negativo: pierde energía y se enfría.
\end{solution}

\begin{solution}{exo:g12:climate-change:2}
Máxima en mayo, mínima en septiembre. En primavera y verano, la vegetación
del hemisferio norte, donde crece la mayoría de las plantas terrestres, capta más
dióxido de carbono por fotosíntesis del que liberan las plantas y los suelos; en otoño
e invierno dominan la descomposición y la respiración, y la concentración vuelve a subir.
\end{solution}

\begin{solution}{exo:g12:climate-change:3}
\emph{Muy probable}: una probabilidad del 90 al 100\,\%; \emph{probable}: del 66 al
100\,\%. $[0.95 \text{ a } 1.20]$ es un intervalo muy probable: hay un 90\,\% de probabilidad de que
el valor verdadero esté entre 0.95 y 1.20.
\end{solution}

\begin{solution}{exo:g12:climate-change:4}
El 91\,\%. El agua se dilata al calentarse: un océano más cálido ocupa más volumen, así que
su superficie sube aunque no se le añada agua (la fusión del hielo continental añade agua
además de esto).
\end{solution}

\begin{solution}{exo:g12:climate-change:5}
Vapor de agua: un aire más cálido contiene más vapor de agua, un gas de efecto invernadero, que se suma
al calentamiento. Nieve y hielo: al fundirse dejan al descubierto superficies más oscuras, que absorben
más luz solar. La respuesta negativa: un planeta más cálido emite más infrarrojo
hacia el espacio, lo que lo devuelve al equilibrio.
\end{solution}

\begin{solution}{exo:g12:climate-change:6}
$427.35 - 315.98 = \qty{111.37}{ppm}$; $111.37/315.98 \approx 0.352$, un
aumento de alrededor del 35\,\%. En 66 años: $111.37/66 \approx
\qty{1.7}{ppm}$ al año de media.
\end{solution}

\begin{solution}{exo:g12:climate-change:7}
$F = 3.93 \times \ln(427.35/278)/\ln 2 = 3.93 \times 0.4300/0.6931 \approx
\qty{2.44}{W/m^2}$.
\end{solution}

\begin{solution}{exo:g12:climate-change:8}
$3.7 \times 30 = \qty{111}{mm}$, unos \qty{11}{cm}. $3.7/1.3 \approx 2.8$:
casi tres veces más rápido.
\end{solution}

\begin{solution}{exo:g12:climate-change:9}
$\Delta T = 2.72/1.16 \approx \qty{2.3}{\degreeCelsius}$. El clima no está
en equilibrio: el océano, que absorbe alrededor del 90\,\% de la energía adicional,
se calienta despacio, así que la superficie aún no se ha puesto a la altura del \omterm{def:g12:climate-change:forcing}{forzamiento} (el
sistema todavía gana unos \qty{0.8}{W/m^2}).
\end{solution}

\begin{solution}{exo:g12:climate-change:10}
El Sol y los volcanes explican como mucho \qty{0.1}{\degreeCelsius} de un
calentamiento de \qty{1.06}{\degreeCelsius}: no pueden ser su causa. La otra
causa natural es la variabilidad interna del clima (sus fluctuaciones
naturales, como los intercambios de calor entre el océano y la atmósfera); la
evaluación sitúa su contribución entre $-0.2$ y $+\qty{0.2}{\degreeCelsius}$,
también demasiado pequeña.
\end{solution}

\begin{solution}{exo:g12:climate-change:11}
Muy probablemente por debajo de \qty{2}{\degreeCelsius}: solo SSP1-1.9 (intervalo muy probable
de 1.0 a 1.8). Muy probablemente por encima: SSP2-4.5 (2.1 a 3.5), SSP3-7.0 (2.8 a 4.6)
y SSP5-8.5 (3.3 a 5.7). Para SSP1-2.6 (1.3 a 2.4), ninguna de las dos cosas es muy probable.
\end{solution}

\begin{solution}{exo:g12:climate-change:12}
El \omterm{def:g12:climate-change:forcing}{forzamiento} de los aerosoles es $-0.22 - 0.84 = \qty{-1.06}{W/m^2}$; sin él, el
total pasa a ser $2.72 + 1.06 = \qty{3.78}{W/m^2}$. El calentamiento de equilibrio
aumenta en $1.06/1.16 \approx \qty{0.9}{\degreeCelsius}$.
\end{solution}

\begin{solution}{exo:g12:climate-change:13}
Como $\ln(2C) - \ln C = \ln 2$ sea cual sea $C$,
\[
  F(556) - F(278) = \frac{3.93\,\ln 2}{\ln 2} = \qty{3.93}{W/m^2}, \qquad
  F(1112) - F(556) = \frac{3.93\,\ln 2}{\ln 2} = \qty{3.93}{W/m^2}:
\]
el mismo \omterm{def:g12:climate-change:forcing}{forzamiento} para cada duplicación. Total a \qty{1112}{ppm}: $2 \times 3.93 =
\qty{7.86}{W/m^2}$.
\end{solution}

\begin{solution}{exo:g12:climate-change:14}
$\dfrac{\dd F}{\dd C} = \dfrac{F_{2\times}}{C \ln 2}$. A \qty{280}{ppm}:
$3.93/(280 \times 0.6931) \approx \qty{0.0202}{W/m^2}$ por ppm; a
\qty{420}{ppm}: $\approx \qty{0.0135}{W/m^2}$ por ppm. Cada ppm adicional añade
menos \omterm{def:g12:climate-change:forcing}{forzamiento} a medida que crece la concentración (1.5 veces menos a 420 que a 280),
pero la concentración sigue subiendo, así que el \omterm{def:g12:climate-change:forcing}{forzamiento} sigue creciendo.
\end{solution}

\begin{solution}{exo:g12:climate-change:15}
$3.93/3.22 \approx \qty{1.2}{\degreeCelsius}$ sin \omterm{def:g12:climate-change:feedback}{retroalimentaciones};
$3.93/1.16 \approx \qty{3.4}{\degreeCelsius}$ con ellas. Factor de
amplificación: $3.22/1.16 \approx 2.8$.
\end{solution}

\begin{solution}{pb:g12:climate-change:1}
\textbf{1.} $3.37 + 3.51 = \qty{6.88}{ppm}$, unos \qty{7}{ppm}.

\textbf{2.} De la primavera al final del verano, las plantas del hemisferio
norte retiran dióxido de carbono del aire por fotosíntesis más deprisa de lo que
la respiración y la descomposición lo devuelven: la concentración baja hasta septiembre.
En otoño e invierno el balance se invierte, y la concentración vuelve a subir
hasta su máximo de mayo.

\textbf{3.} El hemisferio sur tiene mucha menos tierra emergida, y por tanto mucha menos
vegetación, que el norte; la mayor parte de su superficie es océano.

\textbf{4.} A esa altitud está por encima de la capa de aire contaminado cercana al
suelo; lejos de los continentes y de las ciudades, el aire que muestrea está bien mezclado y
representa a toda la atmósfera, no a una fuente local.

\textbf{5.} $(325.68 - 316.91)/10 \approx \qty{0.88}{ppm}$ al año.

\textbf{6.} $(427.35 - 401.01)/10 \approx \qty{2.63}{ppm}$ al año.

\textbf{7.} $2.634/0.877 \approx 3.0$: el ritmo de crecimiento se ha triplicado.

\textbf{8.} $(427.35 - 278)/278 \approx 0.537$: alrededor del 54\,\%.

\textbf{9.} $427.35 - 300 \approx \qty{127}{ppm}$.

\textbf{10.} $F = 3.93 \times \ln(427.35/278)/\ln 2 \approx
\qty{2.44}{W/m^2}$.

\textbf{11.} $F(316.91) = 3.93 \times \ln(316.91/278)/\ln 2 \approx
\qty{0.74}{W/m^2}$; aumento: $2.44 - 0.74 \approx \qty{1.70}{W/m^2}$.

\textbf{12.} Un \unit{W/m^2} más multiplica la concentración por
$2^{1/3.93} \approx 1.193$: $427.35 \times 1.193 \approx \qty{510}{ppm}$.

\textbf{13.} $(509.8 - 427.35)/2.634 \approx 31$ años.

\textbf{14.} $2 \times 278 = \qty{556}{ppm}$; \omterm{def:g12:climate-change:forcing}{forzamiento}: $F_{2\times} =
\qty{3.93}{W/m^2}$.

\textbf{15.} Por definición de la sensibilidad, \qty{3}{\degreeCelsius}
(probablemente de 2.5 a \qty{4}{\degreeCelsius}).

\textbf{16.} $\lambda = 3.93/3 \approx \qty{1.31}{W.m^{-2}.\degreeCelsius^{-1}}$.
La suma de las \omterm{def:g12:climate-change:feedback}{retroalimentaciones}, \qty{1.16}{W.m^{-2}.\degreeCelsius^{-1}}, da
$3.93/1.16 \approx \qty{3.4}{\degreeCelsius}$, algo más; la mejor estimación
evaluada combina esta física con el calentamiento observado y los climas del
pasado, y ambos valores están dentro del intervalo probable.

\textbf{17.} El océano almacena la mayor parte de la energía adicional y se calienta despacio: la
superficie solo alcanza su calentamiento de equilibrio al cabo de décadas o siglos.

\textbf{18.} Supone una historia de las emisiones futuras (el crecimiento de los
últimos diez años mantenido para siempre): eso es un \omterm{def:g12:climate-change:scenario}{escenario}. Una \omterm{def:g12:climate-change:scenario}{proyección} calcula la
respuesta del clima y del ciclo del carbono a un \omterm{def:g12:climate-change:scenario}{escenario de emisiones}
con modelos físicos, y da intervalos y probabilidades.

\textbf{19.} $(556 - 427.35)/2.634 \approx 48.8$ años después de 2025: hacia
\textbf{2074}.
\end{solution}
